\documentclass[11pt,a4paper]{article}
\usepackage[T1]{fontenc}
\usepackage{lmodern,amsmath,amssymb,booktabs,microtype}
\usepackage[margin=24mm]{geometry}
\usepackage[colorlinks=true,urlcolor=blue,linkcolor=blue]{hyperref}
\newcommand{\EE}{\mathbb E}
\newcommand{\tr}{\operatorname{tr}}
\newcommand{\Rem}{\operatorname{Rem}}
\newcommand{\W}{\mathsf W_2}
\newcommand{\Id}{\mathcal I}
\title{Transporting quadratic products across central charge windows\\Output-deficit norms and a coherent-cross calculation}
\author{Alan Day \and Chengheng\thanks{AI research assistant (OpenAI).}}
\date{Local research note --- 6 November 2026}
\begin{document}\maketitle
\begin{abstract}
An adjacent-charge Hilbert--Schmidt isometry transfers the normalized
quadratic product from half filling to a compact central window.
The coherent-force norm bounds each step; a sixth output moment
controls the accumulated deficit-weighted error. This proves the
central-window Gram on the energy square, input deficit and output
deficit, with an explicit $O(p^{-2/3})$ transport bound for the last
field. A candidate remaining force cross is supported by exact
finite CAR contractions and full-rank sums. Its general-charge
Casimir recoupling step is isolated at the end for further proof.
\end{abstract}
\section{State, normalization and the established result}
Use the full overlapping complex Gaussian $p$-body Hamiltonian,
standard CAR, $\hbar=1$, and normalized maximally mixed trace in
charge $m$. Norms are annealed:
$\|T\|_{2,m}^2=\EE\tr_m(T^\dagger T)$. All sectors share the
same Gaussian couplings and variance $w_N$. Take even $N$ and
\begin{equation}
 \lambda_N=p^2/N\to\lambda>0,\quad m=N/2+xp,\quad |x|\le L,
 \qquad q=e^{-4\lambda}.
\end{equation}
Bounds are uniform for $\lambda_N$ in a positive compact set and
integer $m$ in this window. Let $h_a$ be its rank-$a$ Gaussian part,
$C_m=\EE h^2=w_N\binom mp\binom{N-m+p}p$, and $\W$ denote ordinary
Gaussian second chaos. Put
\begin{align}
 B_m&=\W h^2,& D_m&=\sum_a(p-a)h_a,\
 L_{\rm in,m}&=\W\{h,D_m-\bar u_mh\},&
 \bar u_m&=\sum_a(p-a)\EE h_a^2/C_m,\
 \mathcal S_m&=\sum_{a\le b,r}(a+b-r)\W Z^m_{ab;r}.&&
\end{align}
Here $Z^m_{ab;r}=P_r(h_a\odot h_b)$, where $\odot$ is a square
for $a=b$ and an anticommutator otherwise. We prove
\begin{equation}
 \frac{G(B,L_{\rm in},\mathcal S)}{C_m^2}\longrightarrow
 \begin{pmatrix}1+q&0&b\\0&2\lambda(1+q)&0\\b&0&s\end{pmatrix},
 \quad\begin{cases}b=2\lambda(1-q),\\
 s=4\lambda^2(1+q)+2\lambda(1-q).\end{cases}\label{three}
\end{equation}
The known half-filled Gram is the reference, not a parity identity
assumed at the destination charge.

\newpage
\section{Exact adjacent-charge transport}
Set $n=N-m+1$ and $x_r=r(N+1-r)$. The established removal norm is
\begin{equation}
 \|\Rem T_r\|_{2,m}^2=\ell_{m,r}^2\|T_r\|_{2,m-1}^2,
 \qquad \ell_{m,r}^2=\frac m n(mn-x_r),\quad
 \Rem T=\sum_i c_i^\dagger T c_i.\label{rem}
\end{equation}
Define $\Id_m=\sum_{r\le2p}\ell_{m,r}^{-1}\Rem P_r$.
For large $p$ the entire central path satisfies
$2p\le\min(m-1,N-m)$, so both sectors contain every rank used here.
Equivariance and multiplicity-free neutral operator spaces make
$\Id_m$ an onto isometry on this common subspace. Its inverse is
its normalized-HS adjoint. No assertion about positivity of this
map is required.

The map acts only on the output matrix of a Gaussian coefficient
tensor. The common $p$-body coupling space decomposes into the same
input ranks in all sectors. Thus input-pair blocks and output ranks
are preserved, including their internal multiplicities. The map
commutes with $\W$ and every block weight $f(a,b,r)$.

Write $\widehat B^m_{ab;r}=\W Z^m_{ab;r}/C_m$ and
\begin{align}
 R_m&=C_{m-1}/C_m=\frac{(m-p)(n+p)}{mn},&t_{m,r}&=R_m\ell_{m,r},\
 A_m&=\Rem(H_{m-1}^2)+(2p-m)H_m^2,&O_m&=\W A_m.
\end{align}
The same-$J$ product identity is exactly
\begin{equation}
 t_{m,r}\Id_m\widehat B^{m-1}_{ab;r}-(m-2p)\widehat B^m_{ab;r}
       =O^m_{ab;r}/C_m.\label{identity}
\end{equation}
In the stated window, uniformly for $r\le2p$,
\begin{equation}
 t_{m,r}\ge cN,\qquad |t_{m,r}-(m-r)|\le K,\qquad
 |t_{m,r}-(m-2p)|\le K(1+u+v+S),\label{scalar}
\end{equation}
where $u=p-a,v=p-b,S=a+b-r$. These follow by squaring $t$:
all scaled denominators are bounded away from zero, and
$t^2-(m-r)^2=O(N)$. More explicitly, for $r=\rho p$,
\begin{align}
 t_{m,r}-(m-r)&\longrightarrow2\lambda[x(2-\rho)-1],\
 p^2(R_m-1)&\longrightarrow4\lambda^2(2x-1).
\end{align}
The first limit is a scalar check for bounded $\rho\in[0,2]$;
the proof below uses the uniform inequality \eqref{scalar}.

\newpage
\section{Telescoping and the output tail}
Two inputs from the preceding central-window notes are used:
\begin{align}
 \sum_{a\le b,r}(1+u+v)^2\|O^m_{ab;r}/C_m\|_2^2&\le K,\label{forcebound}\
 \sum_{a\le b,r}S^6\|\widehat B^m_{ab;r}\|_2^2&\le K.\label{six}
\end{align}
There are also uniform factorial tails in $u,v$. For clarity,
\eqref{six} follows from the positive ordered-product law $\pi$;
the symmetrized same-$J$ block is bounded by twice the ordered
law after the ordered-pair convention is included. It does not
require a sign formula for the exchanged contraction. The same
bound holds with each fixed polynomial in $u+v$ inserted, by the
conditional uniform sixth moment and the factorial input tail.

Solving \eqref{identity}, and summing squared orthogonal block norms,
now gives
\begin{equation}
 \left\|\Id_m(B_{m-1}/C_{m-1})-B_m/C_m\right\|_2\le K/N.
 \label{step}
\end{equation}
Indeed the numerator is $O_m/C_m+(m-2p-t_{m,r})\widehat B^m$
in each block; \eqref{forcebound}, \eqref{scalar} and the product
second moment bound it uniformly. The same argument with weight
$u+v$ uses \eqref{forcebound} and the weighted product moments.
Composing the isometries (their inverses below half filling) over
at most $Lp$ steps proves an unweighted error $Kp/N=O(p^{-1})$.

Let $U_{m\leftarrow0}$ denote this composition from $m_0=N/2$.
All steps preserve $a,b,r$, so the cutoff $P_{S\le R}$ is the same
coefficient-space projector throughout. Telescope first and then
multiply by $S$ on this low block. For the two endpoint tails,
\eqref{six} gives
\begin{equation}
 \|S\mathbf1_{S>R}\widehat B^m\|_2^2\le K/R^4.
\end{equation}
Consequently
\begin{equation}
 \left\|U_{m\leftarrow0}\frac{\mathcal S_{m_0}}{C_{m_0}}
                   -\frac{\mathcal S_m}{C_m}\right\|_2
 \le K\left(\frac R p+\frac1{R^2}\right)
 =O(p^{-2/3})\quad(R=p^{1/3}).\label{transport}
\end{equation}
Only the two endpoint tails are counted; summing a tail bound at
every intermediate charge would give a different, inadequate estimate.

The input-weighted version of \eqref{step}, together with
$\bar u_m\to\lambda$ uniformly, also transports $L_{\rm in}/C_m$
with an error tending to zero. Thus all three fields in \eqref{three}
converge in the same coefficient Hilbert space to their half-filled
references. Uniform boundedness and Cauchy--Schwarz transfer every
Gram entry, proving \eqref{three}. The displayed rate concerns
transport only, not the rate of the reference Gram's DS limit.

\newpage
\section{A candidate remaining coherent cross}
The preceding proof establishes the output self norm. To add $O_m$
to the Gram, consider $X=x_{a+b}-x_r$. Since
\begin{equation}
 X=S[N+1-2(a+b)+S],\qquad
 \|\mathcal X_m/N-\mathcal S_m\|_2/C_m=O(p^{-1}),\label{proxy}
\end{equation}
the uniform $O_m/C_m$ bound permits replacing the deficit by this
Casimir proxy in a cross. Define the neutral Racah kernel $\chi$
with lower parameters $(-m,m-N,-N-1)$, and
\begin{equation}
 \xi={}_4F_3\!\left[\begin{matrix}-a,a-N-1,-b,b-N-1\\
 1-m,m-N,-N-1\end{matrix};1\right],\qquad D=\xi-\chi.
\end{equation}
Both kernels have the same displayed upper parameters. Put
$g=(m-p)^2/m$, $\eta=p^2/m$ and
\begin{equation}
 C_f=\frac{m(N-m)(N+1)}N,\quad
 C_t=\frac{(m-1)(N-m)(N+1)}{N-1},\quad
 U_{ab}=-2ab(1+\chi)+x_ax_b/C_f.
\end{equation}
The neutral proxy cross per ordered pair is $U_{ab}/N$.
For the coherent product with $\Rem(H_{m-1}^2)$, the candidate
direct and exchanged Casimir moments, divided by $C_aC_b$, are
\begin{equation}
 g(x_a+x_b-x_ax_b/C_t),\qquad
 g[(x_a+x_b)\xi-mD].\label{candidatelemma}
\end{equation}
They pass the independent pure-rank CAR checks described on page 6.
If \eqref{candidatelemma} holds generally, algebra gives the stable
same-$J$ formula
\begin{equation}
 z_{ab}=\frac1N\left\{\eta U_{ab}
 +g\left[\frac{x_ax_b}{m(m-1)(N+1)}+(m-2ab)D\right]\right\}.
 \label{candidate}
\end{equation}
It combines the large terms before evaluation, using
$C_t^{-1}-C_f^{-1}=[m(m-1)(N+1)]^{-1}$.

For fixed input deficits, terminating-series bounds give
$\chi\to q$ and $mD\to-4\lambda q$. Each term is bounded by
$K^j/j!$, including one extra $j$ moment. Thus \eqref{candidate}
is uniformly bounded over all $a,b\le p$ and tends to $s$.
The factorial input tail justifies the full rank sum. Combined
with \eqref{proxy}, this would prove
$\langle\mathcal S_m,O_m\rangle/C_m^2\to s$.
This last inference is conditional on the general-charge
recoupling identity \eqref{candidatelemma}, not on numerical convergence.

\newpage
\section{The precise analytic step still needed}
The charged covariance is diagonal with terminating kernel
\begin{equation}
 T_a(t)={}_4F_3\!\left[\begin{matrix}-a,a-N-1,-t,t-N\\
 1-m,m-N,-N-1\end{matrix};1\right].
\end{equation}
The established two-component decomposition gives
\begin{align}
 \xi&=T_a(b)-\frac b{N+1-2b}[T_a(b-1)-T_a(b)],\
 \frac{x_b}{N+1-2b}[T_a(b-1)-T_a(b)]&=-mD.\label{termwise}
\end{align}
The second equality follows termwise: if $c_j$ is a summand of
$\chi$, the corresponding summand of $\xi$ is $mc_j/(m-j)$.
We have additionally checked 624 finite rational instances.

What remains is to identify the exchanged gradient contraction
with the left side of \eqref{termwise} times $g$. With
$E_{ij}=c_i^\dagger c_j$ and
$\mathcal C=\frac12\sum_{ij}\operatorname{ad}_{E_{ij}}
\operatorname{ad}_{E_{ji}}$, the exact product rule is
\begin{equation}
 \mathcal C(UV)=(x_a+x_b)UV+
          \sum_{ij}[E_{ij},U][E_{ji},V].
\end{equation}
The needed evaluation uses the two charged components of $Vc_i$
and the CAR commutators $[E_{ij},c_k]=-\delta_{ik}c_j$.
It must track the rectangular trace normalization and both
component weights. The termwise identity alone does not evaluate
this operator contraction. The half-filled alternating-completeness
proof in the earlier note does not supply that general-charge step.

For a concrete continuation, set $z=1+q$, $l=2\lambda(1+q)$,
$a_*=4\lambda+4\lambda^2(1+q)$. The established subblocks
\eqref{three} and the preceding force note leave exactly one unknown:
\begin{equation}
 C_m^{-2}G(B,L_{\rm in},\mathcal S,O)\ \longrightarrow\
 \begin{pmatrix}z&0&b&b\\0&l&0&0\\b&0&s&c\\b&0&c&a_*\end{pmatrix}.
\end{equation}
Existence and value of $c$ are the remaining cross problem.
The candidate is $c=s$. Under that condition, writing
\begin{equation}
 v=s-b^2/z=2\lambda(1-q)+\frac{16\lambda^2q}{1+q},\qquad
 \Delta=a_*-s=2\lambda(1+q),
\end{equation}
the residual of $\mathcal S/C_m$ after projection onto
$\operatorname{span}(B,L_{\rm in},O)$ has variance
$v\Delta/(v+\Delta)>0$. This is a conditional Schur complement.
Its central-window status must remain distinct from the already
established half-filled residual.

\newpage
\section{Reproducible checks and continuation}
The new files in \texttt{results/susskind-memory-central-cross}
record the following independent scopes:
\begin{itemize}
 \item \texttt{transport.json}: exact symbolic scalar drift and
 variance ratio; complete rational isometry matrices on ranks
 $0,1,2$ for $(N,p,m)=(6,1,3)$, dimensions $1,35,189$.
 Full 1,296 ordered Gaussian coordinate pairs verify the normalized
 transport identity with maximal error $3.47\,10^{-17}$; the
 square-root isometry norm comparison differs by $2.23\,10^{-16}$.
 \item \texttt{charged-casimir.json}: 15 pure-rank pairs in four
 small systems, all 60 direct/exchanged zeroth/first Casimir
 identities exact. Across these blocks 66,531 ordered Gaussian
 coordinate pairs are contracted, without a parity substitution.
 \item \texttt{car-exploration.json}: independent full $4\times4$
 coefficient Grams on $(B,\mathcal S,\mathcal X/N,O)$, totaling
 31,466 symmetric coordinate pairs. All four differences between
 the actual proxy--force cross and the candidate are exactly zero.
 Its original exploratory status is retained.
 \item \texttt{cross-limit.json}: 624 rational contiguous checks,
 four finite all-rank exact comparisons, and 36 central cases at
 $\lambda=0.1,0.25,0.5$, $p=16,32,64,128$, $x=-2,0,2$.
 Every input rank and terminating-series term is retained.
 Evaluations at 80 and 120 digits agree at the stored 65 significant
 digits. The three columns are $\mathcal X/N$ crossed with $B$,
 $L_{\rm in}$ and the candidate $O$ expression, respectively.
\end{itemize}
At $\lambda=0.5,p=128$, the candidate force cross divided by
$C_m^2$ is approximately $1.97682736$, $1.94517151$, $1.91541182$
for $x=-2,0,2$; its proposed limit is $2$. These finite-size
deviations are much larger than precision-refinement differences.
Neither floating refinement nor exact small-system checks prove
the unresolved recoupling lemma.

The transport theorem uses even $N$ and a fixed compact central
window. Odd-$N$ references, growing charge windows and full charge
sums require additional work. The next immediate task is the
gradient contraction on page 5, followed by energy insertions,
the raw memory spectrum and finite-memory remainder. The
space/mass dictionary for Susskind's bulk mesons remains downstream.

\begin{thebibliography}{9}
\bibitem{TT} J. Tall and S. Tomsovic,
\emph{Double-scaled bosonic and fermionic embedded ensembles,
complex SYK, and the dual Hilbert space},
\href{https://arxiv.org/html/2604.14522v2}{arXiv:2604.14522v2},
Sec. 2.1, Eqs. (2.11)--(2.19). The irreducible covariance and
fixed-charge framework are prior results; particle--hole symmetry
does not imply a same-sector product parity away from half filling.
\bibitem{local} A. Day and Chengheng, local project notes:
\emph{Parity-weighted rank deficits},
\emph{Central charge windows}, and
\emph{Coherent memory norms on central charge windows}
(6 November 2026). These provide the reference Gram, positive
sixth moment and coherent-force bound used in the transport proof.
\end{thebibliography}
\end{document}
